\documentclass[12pt]{article}

\input{../../_assets/preambulo.tex}


\begin{document}

    % 1. Foto de fondo
    % 2. Título
    % 3. Encabezado Izquierdo
    % 4. Color de fondo
    % 5. Coord x del titulo
    % 6. Coord y del titulo
    % 7. Fecha

    
    \input{../../_assets/portada}
    \portadaExamen{ffccA4.jpg}{Topología I\\Examen XIV}{Topología I. Examen XIV}{MidnightBlue}{-8}{28}{2026}{Roxana Acedo Parra}

    \begin{description}
        \item[Asignatura] Topología I.
        \item[Curso Académico] 2025/2026.
        \item[Grado] Grado en Matemáticas.
        \item[Grupo] Único.
        \item[Profesor] Antonio Alarcón López.
        \item[Descripción] Convocatoria Ordinaria.
        \item[Fecha] 14 de enero de 2026.
        \item[Duración] 180 minutos.
    
    \end{description}
    \newpage


    % ------------------------------------
    
    \begin{ejercicio}
        En $\bb{Z}$ se considera la familia:
        \[
        \cc{B} = \left\{\left\{2m - 1, 2m, 2m + 1\right\}, m \in \bb{Z}\right\} \cup \left\{\left\{2m + 1\right\}, m \in \bb{Z}\right\}.
        \]
        \begin{enumerate}[label=\alph*)]
            \item Prueba que $\cc{B}$ es base de una topología $\cc{T}_K$ en $\bb{Z}$. Al espacio topológico $(\bb{Z}, \cc{T}_K)$ se le denomina recta de Khalimsky.
            \item Sea $A = \{-3, -2, 1, 2, 3\} \subset \bb{Z}$. Calcula el interior y la clausura de $A$.
            \item Dado $(X, \cc{T})$ espacio topológico, se dice que $(X, \cc{T})$ es $T_0$, si para cada par de puntos $x, y \in X$, $x \neq y$ existe $V_x$ entorno de $x$ tal que $y \notin V_x$ o existe $V_y$ entorno de $y$ tal que $x \notin V_y$. Prueba que $(\bb{Z}, \cc{T}_K)$ es $T_0$. ¿Es $T_1$?
            \item Estudia la conexión y la compacidad de $(\bb{Z}, \cc{T}_K)$.
            \item Comprueba que $f : (\bb{Z}, \cc{T}_K) \to (\bb{Z}, \cc{T}_K)$ es continua si y solo si:
            \begin{description}
                \item[e1)] Si $x$ es par y $f(x)$ es impar entonces $f(x \pm 1) = f(x)$.
                \item[e2)] $|f(x \pm 1) - f(x)| \leq 1$, para todo $x \in \bb{Z}$.
            \end{description}
            \item Consideremos en $\bb{R}$ la relación de equivalencia
            
            $$x R x' \sii x = x' \quad \text{o} \quad x, x' \in \left[2m - \nicefrac{1}{2}, 2m + \nicefrac{1}{2}\right] \quad \text{o} $$ $$ \quad x, x' \in \left](2m+1) - \nicefrac{1}{2}, (2m+1) + \nicefrac{1}{2}\right[, \quad m \in \bb{Z}.$$
            
            Demuestra que $(\bb{R}/R, \cc{T}_u/R)$ es homeomorfo a $(\bb{Z}, \cc{T}_K)$.
        \end{enumerate}
    \end{ejercicio}

    \begin{ejercicio}
        Enuncia el Lema del Tubo y el Teorema de Tychonoff. Demuestra el Teorema de Tychonoff.
    \end{ejercicio}

    \begin{ejercicio}
        Una aplicación $f : (X, \cc{T}) \to (Y, \cc{T}')$ entre espacios topológicos se dice propia si para todo $K$ compacto de $(Y, \cc{T}')$ se tiene que $f^{-1}(K)$ es un compacto de $(X, \cc{T})$.
        \begin{enumerate}[label=\alph*)]
            \item Sea $(X, \cc{T})$ un espacio topológico $T_2$ e $(Y, \cc{T}')$ un espacio topológico compacto. Prueba que si $f : (X, \cc{T}) \to (Y, \cc{T}')$ es propia entonces $f$ es continua. ¿Es cierto el recíproco?
            \item Prueba que toda aplicación continua $f : (\bb{S}^1, \cc{T}) \to (\bb{S}^1, \cc{T})$ es propia, donde $\cc{T}$ es la topología inducida de la topología usual de $\bb{R}^2$.
        \end{enumerate}
    \end{ejercicio}

    \newpage
    \setcounter{ejercicio}{0} % Reiniciar contador de ejercicios
    % \noindent
    % \textbf{Solución.}

        \begin{ejercicio}
        En $\bb{Z}$ se considera la familia:
        \[
        \cc{B} = \left\{\left\{2m - 1, 2m, 2m + 1\right\}, m \in \bb{Z}\right\} \cup \left\{\left\{2m + 1\right\}, m \in \bb{Z}\right\}.
        \]
        \begin{enumerate}[label=\alph*)]
            \item Prueba que $\cc{B}$ es base de una topología $\cc{T}_K$ en $\bb{Z}$. Al espacio topológico $(\bb{Z}, \cc{T}_K)$ se le denomina recta de Khalimsky.

                Veamos que se verifican las propiedades B1 y B2. En primer lugar es claro que
                \[
                \mathbb{Z} = \bigcup_{m \in \mathbb{Z}} \{2m-1, 2m, 2m+1\},
                \]
                por lo que $\mathbb{Z} = \bigcup_{B \in \mathcal{B}} B$.
                
                En segundo lugar estudiamos los casos en los que dos elementos distintos de $\mathcal{B}$ tienen intersección no vacía:
                
                \textbf{Caso 1:} $\{2m-1, 2m, 2m+1\} \cap \{2m+1\} = \{2m+1\} \in \mathcal{B}$ y $\{2m-1, 2m, 2m+1\} \cap \{2m-1\} = \{2m-1\} \in \mathcal{B}$.
                
                \textbf{Caso 2:} $\{2m-1, 2m, 2m+1\} \cap \{2m+1, 2m+2, 2m+3\} = \{2m+1\} \in \mathcal{B}$ y $\{2m-1, 2m, 2m+1\} \cap \{2m-3, 2m-2, 2m-1\} = \{2m-1\} \in \mathcal{B}$.
                
                Por lo tanto, si $B_1, B_2 \in \mathcal{B}$ con $B_1 \cap B_2 \neq \emptyset$ entonces $B_1 \cap B_2 \in \mathcal{B}$ y se cumple trivialmente B2.
                
                Observemos que una base de entornos de la topología $\mathcal{T}_K$ es: $\mathcal{B}_x = \{\{x\}\}$ si $x$ es impar, $\mathcal{B}_x = \{\{x-1, x, x+1\}\}$ si $x$ es par.
        
            \item Sea $A = \{-3, -2, 1, 2, 3\} \subset \bb{Z}$. Calcula el interior y la clausura de $A$.
            
                $\text{int}(A) = \{-3, 1, 2, 3\}$ ya que para esos puntos tenemos los entornos abiertos $\{-3\}$ y $\{1, 2, 3\}$ contenidos en $A$, y $\{-3, -2, -1\}$ no está contenido en $A$, por lo que $A$ no es entorno de $-2$.
                
                $\text{cl}(A) = \{-4, -3, -2, 0, 1, 2, 3, 4\}$ ya que los entornos de la base de entornos de esos puntos cortan al conjunto $A$, y si $p \in \mathbb{Z}$ no está en $\{-4, -3, -2, 0, 1, 2, 3, 4\}$ podemos encontrar un entorno de $p$ que no corta a $A$.

            \item Dado $(X, \cc{T})$ espacio topológico, se dice que $(X, \cc{T})$ es $T_0$, si para cada par de puntos $x, y \in X$, $x \neq y$ existe $V_x$ entorno de $x$ tal que $y \notin V_x$ o existe $V_y$ entorno de $y$ tal que $x \notin V_y$. Prueba que $(\bb{Z}, \cc{T}_K)$ es $T_0$. ¿Es $T_1$?

                Veamos que es $T_0$. En primer lugar consideramos $2m+1$ para $m \in \mathbb{Z}$ y $p \in \mathbb{Z}$ con $p \neq 2m+1$. Es claro que $\{2m+1\}$ es un entorno de $2m+1$ que no contiene a $p$. En segundo lugar consideramos $2m, 2m' \in \mathbb{Z}$ con $m < m'$. En ese caso tenemos $\{2m-1, 2m, 2m+1\}$ un abierto de $2m$ que no contiene a $2m'$.
            
                La recta de Khalimsky no es $T_1$ ya que por ejemplo, $1$ y $2$ no verifican la definición de $T_1$. No existe ningún entorno de $2$ que no contenga a $1$.
                
            \item Estudia la conexión y la compacidad de $(\bb{Z}, \cc{T}_K)$.
                    
                Los abiertos básicos de la forma $\{2m-1, 2m, 2m+1\}$ son conexos ya que son los abiertos más pequeños que contienen a $2m$ (por lo que la única separación que admite dicho conjunto es la trivial). Además podemos escribir
                \[
                \mathbb{Z} = \bigcup_{m \in \mathbb{Z}} \{2m-1, 2m, 2m+1\}.
                \]
                Observemos que esto es una unión de conexos tal que cada uno de ellos corta al siguiente. Por tanto $\mathbb{Z}$ es conexo.
                
                Por otra parte es sencillo comprobar que la recta de Khalimsky no es compacta. Efectivamente, podemos considerar el recubrimiento de $\mathbb{Z}$,
                \[
                \mathcal{A} = \{\{2m-1, 2m, 2m+1\} : m \in \mathbb{Z}\}.
                \]
                De este recubrimiento no podríamos sacar un subrecubrimiento finito porque cualquier subrecubrimiento finito sólo contendría un número finito de números enteros.
            
            \item Comprueba que $f : (\bb{Z}, \cc{T}_K) \to (\bb{Z}, \cc{T}_K)$ es continua si y solo si:
            \begin{description}
                \item[e1)] Si $x$ es par y $f(x)$ es impar entonces $f(x \pm 1) = f(x)$.
                \item[e2)] $|f(x \pm 1) - f(x)| \leq 1$, para todo $x \in \bb{Z}$.
            \end{description}

                 Recordemos que $f$ es continua en $x$ si y solo si para toda $B' \in \mathcal{B}_{f(x)}$ existe $B \in \mathcal{B}_x$ tal que $f(B) \subset B'$.
            
                Supongamos que $f$ es continua y sea $x$ un número par. Supongamos en primer lugar que $f(x)$ es impar. Entonces tenemos que $B' = \{f(x)\}$ y $B = \{x-1, x, x+1\}$ y $f$ es continua en $x$ si y solo si $f(\{x-1, x, x+1\}) \subset \{f(x)\}$, es decir si y solo si $f(x\pm 1) = f(x)$.
                
                Supongamos ahora que $f(x)$ es par. Entonces $B' = \{f(x)-1, f(x), f(x)+1\}$ y $B = \{x-1, x, x+1\}$ y $f$ es continua en $x$ si y solo si $\{f(x\pm 1), f(x)\} \subset \{f(x)\pm 1, f(x)\}$.
                
                Observemos que esto y lo anterior implican e2) en el caso $x$ par. Si $x$ es impar para obtener e2) basta aplicar e2) a $x+1$ y $x-1$ que son pares.
                
                Recíprocamente, supongamos que se verifica e1) y e2). Por lo visto en la implicación anterior se tendría que $f$ es continua en los números pares. Notemos que cualquier $f$ es continua en $x$ si $x$ es un número impar ya que en este caso $B = \{x\}$ siempre verifica $f(B) \subset B'$ para cualquier $B'$.
            
            \item Consideremos en $\bb{R}$ la relación de equivalencia
            
            $$x R x' \sii x = x' \quad \text{o} \quad x, x' \in \left[2m - \nicefrac{1}{2}, 2m + \nicefrac{1}{2}\right] \quad \text{o} $$ $$ \quad x, x' \in \left](2m+1) - \nicefrac{1}{2}, (2m+1) + \nicefrac{1}{2}\right[, \quad m \in \bb{Z}.$$
            
            Demuestra que $(\bb{R}/R, \cc{T}_u/R)$ es homeomorfo a $(\bb{Z}, \cc{T}_K)$.

                Consideremos la aplicación $f: (\mathbb{R}, \mathcal{T}_u) \rightarrow (\mathbb{Z}, \mathcal{T}_K)$ dada por
                \[
                f(x) = \begin{cases}
                2m & \text{si } x \in [2m - \frac{1}{2}, 2m + \frac{1}{2}] \\
                2m+1 & \text{si } x \in ](2m+1) - \frac{1}{2}, (2m+1) + \frac{1}{2}[
                \end{cases}
                \]
                Veamos que $f$ es una identificación. Claramente $f$ es sobreyectiva. Veamos que $\mathcal{T}_K = (\mathcal{T}_u)_f$.
                
                Observemos en primer lugar que $f^{-1}(\{2m+1\}) = ](2m+1) - \frac{1}{2}, (2m+1) + \frac{1}{2}[$ y $f^{-1}(\{2m-1, 2m, 2m+1\}) = ](2m-1) - \frac{1}{2}, (2m+1) + \frac{1}{2}[$ son abiertos de $(\mathbb{R}, \mathcal{T}_u)$. Por tanto $f$ es continua y esto implica que $\mathcal{T}_K \subset (\mathcal{T}_u)_f$.
                
                Falta ver que $(\mathcal{T}_u)_f \subset \mathcal{T}_K$. Sea $O \in (\mathcal{T}_u)_f$, entonces sabemos que $f^{-1}(O) \in \mathcal{T}_u$ por ser $(\mathcal{T}_u)_f$ la topología final para la aplicación $f$. Veamos que $O \in \mathcal{T}_K$ comprobando que es entorno de todos sus puntos. Si $2m+1 \in O$ para algún $m \in \mathbb{Z}$, entonces $\{2m+1\} \subset O$ y $O$ es entorno de $2m+1$. Consideremos ahora $2m \in O$ para $m \in \mathbb{Z}$ entonces $[2m - \frac{1}{2}, 2m + \frac{1}{2}] \subset f^{-1}(O)$ y $2m \pm \frac{1}{2} \in f^{-1}(O) \in \mathcal{T}_u$. Por tanto, existe $0 < \epsilon < 1/2$ tal que $(2m - \frac{1}{2} - \epsilon, 2m + \frac{1}{2} + \epsilon) \subset f^{-1}(O)$. Luego $\{2m-1, 2m, 2m+1\} = f((2m - \frac{1}{2} - \epsilon, 2m + \frac{1}{2} + \epsilon)) \subset O$, y $O$ es entorno de $2m$. Concluimos que $O \in \mathcal{T}_K$.
                
                Es claro que $R = R_f$ por definición de la aplicación $f$ y por tanto la aplicación $\tilde{f}: (\mathbb{R}/R, \mathcal{T}_u/R) \rightarrow (\mathbb{Z}, \mathcal{T}_K)$ es un homeomorfismo.
        \end{enumerate}
    \end{ejercicio}

    \begin{ejercicio}
        Enuncia el Lema del Tubo y el Teorema de Tychonoff. Demuestra el Teorema de Tychonoff. \\
        Teoría dada en clase.
    \end{ejercicio}

    \begin{ejercicio}
        Una aplicación $f : (X, \cc{T}) \to (Y, \cc{T}')$ entre espacios topológicos se dice propia si para todo $K$ compacto de $(Y, \cc{T}')$ se tiene que $f^{-1}(K)$ es un compacto de $(X, \cc{T})$.
        \begin{enumerate}[label=\alph*)]
            \item Sea $(X, \cc{T})$ un espacio topológico $T_2$ e $(Y, \cc{T}')$ un espacio topológico compacto. Prueba que si $f : (X, \cc{T}) \to (Y, \cc{T}')$ es propia entonces $f$ es continua. ¿Es cierto el recíproco?
            
                Comprobaremos que si $C$ es un cerrado de $(Y, \mathcal{T}')$ entonces $f^{-1}(C)$ es un cerrado de $(X, \mathcal{T})$. Observemos que $C$ es compacto por ser un cerrado de un espacio topológico compacto. Como $f$ es propia tenemos entonces que $f^{-1}(C)$ es un compacto de $(X, \mathcal{T})$ que es $T_2$ y por tanto $f^{-1}(C)$ es un cerrado de $(X, \mathcal{T})$. El recíproco no es cierto ya que la aplicación inclusión $i: ((1, 2), (\mathcal{T}_u)_{(1,2)}) \rightarrow ([0, 3], (\mathcal{T}_u)_{[0,3]})$ es continua pero no es propia.
                
            \item Prueba que toda aplicación continua $f : (\bb{S}^1, \cc{T}) \to (\bb{S}^1, \cc{T})$ es propia, donde $\cc{T}$ es la topología inducida de la topología usual de $\bb{R}^2$.
            
                Consideremos $K$ compacto de $\mathbb{S}^1$. Como $\mathbb{S}^1$ es $T_2$ tenemos que $K$ es cerrado. Usando que $f$ es continua tenemos entonces que $f^{-1}(K)$ es un cerrado de $\mathbb{S}^1$ que es compacto y por tanto $f^{-1}(K)$ es compacto.
        \end{enumerate}
    \end{ejercicio}

\end{document}